\scanpage{3}
% L1
\deleted{Toute} \unsure{Tout|Toute} \deleted{représentation}\inserted{épi.} du gpe \deleted{\unsure{$L$}}
% L2
\[L_{\deleted{2}3}=\{\rho_0,\rho_1,\rho_2\mid\rho_0\rho_1\rho_2=1\}\]
% L3
dans un gpe fini $\Gamma$ donne une carte finie
% L4
\deleted{homogène} $3$-pondérée homogène (ou régulière,
% L5
ou galoisienne) -- avec l'étude arithmétique
% L6
\unsure{à la} clef. Donc $\forall$ hom de $L_2$ dans un
% L7
$GP(1,k)$, \deleted{donne} à image finie, donne une telle
% L8
carte -- or les représentations en question
% L9
ont une interprétation géométrique
% L10
élémentaire, en termes d'hyperpolygônes
% L11
réguliers. Donc on a une étude
% L12
arithmétique \unsure{directe|de tte} de ceux-ci -- en
% L13
commençant par ceux qui sont
% L14
définis par des systèmes de \unsure{Schwartz|Schwarz}
% L15
\deleted{hyper} pythagoriciens. \unsure{Ceux-ci|Ceux qui} définissent
% L16
44 types de \unsure{'}hyperpolygônes -- \unsure{stables|obtenus} par
% L17
\deleted{symétrie}\inserted{$\unsure{\mathfrak S_3|\mathcal S_3}$} $\times$ ``opérations élémentaires'' --
% L18
et il y a un \deleted{gros}\inserted{(plus petit \unsure{(et fort gros!)|c'est fort gros!} )} quo. \deleted{\unsure{très}} fini de
% L19
$L_2$ qui les coiffe tous.

% L20
\deleted{D'autres exemples sont liés aux}\inserted{Parmi ceux-ci, on a \deleted{ceux} correspondant aux}
% L21
systèmes de Schwartz \uline{canoniques}
% L22
des $n$-gones pythagoriciens de $GP(1)$,
% L23
ou encore, aux bipyramides régulières habituelles
% L24
-- plus leurs transformés par $\mathfrak S_3$.

% margin note written sideways
% L25
\margin{Il convient d'y ajouter ceux qui sont ``dégénérés'', i.e. qui correspondent de trois façons différentes aux représentations de $\mathbb D_\infty\times\mathbb Z_{/2}$, et enfin les hyperdégénérés (représentations \deleted{\unsure{presque}} abéliennes de $G_2$) et \unsure{stabiliser|stabilisés} par $\mathfrak S_3$ et opérations élémentaires\ldots.}
% archive sheet number: 2
